\section*{Conference Papers \& Oral Presentations}
Presenters are indicated by asterisks.\\ 

\normalsize

\begin{etaremune}[topsep=0pt,itemsep=0pt,partopsep=0pt,parsep=0pt]
\item Basu*, S., DeMarco, A. W., and He, P. (2020), ``On the outer length scales of optical turbulence'', Propagation through and Characterization of Atmospheric and Oceanic Phenomena (pcAOP), OSA, online presentation. 
\item Basu, S. (2019). ``Mesoscale modelling of optical turbulence in the atmosphere: quantifying the impact of ultra-high vertical resolution'', Communications and Observations through Atmospheric Turbulence (COAT), Onera, France. 
\item Basu, S. (2019). ``Multiscale modelling of atmospheric turbulence'', workshop on Non-Kolmogorov Turbulence and Associated Phenomena, Fraunhofer IOSB, Ettlingen, Germany.
\item Basu, S. (2019). ``Utilizing the ascent rates of weather balloons to estimate optical turbulence ($C_n^2$) profiles'', Propagation through and Characterization of Atmospheric and Oceanic Phenomena (pcAOP), OSA, Munich, Germany. 
\item Basu*, S., Bose-Pillai, S. R., Fiorino, S. T., and McCrae, J. E. (2018) ``Evaluating a coupled mesoscale modeling and ray tracing framework over an urban area'', Propagation through and Characterization of Atmospheric and Oceanic Phenomena (pcAOP), OSA, Orlando, FL, USA. 
\item Kulikov*, V., Basu, S., and Vorontsov, M. (2017). ``Simulation of laser beam propagation based on mesoscale modeling of optical turbulence and refractivity'', 26-29 June, Propagation through and Characterization of Atmospheric and Oceanic Phenomena (pcAOP), OSA, San Francisco, CA, USA.
\item Basu, S. (2017). ``Revisiting the effects of spatial resolution on large-eddy simulations of stable boundary layers'', 27-31 March, Stable Boundary Layer Workshop, Delft, The Netherlands.  
\item Basu, S. (2017). ``On the usage of M-O similarity theory in very high-resolution LES'', 27-31 March, Stable Boundary Layer Workshop, Delft, The Netherlands. 
\item Kiliyanpilakkil*, V. P., and Basu, S. (2017). ``Understanding the existence of quasi-universal scaling characteristics of the lower atmospheric boundary layer wind speeds in the mesoscale regime'', 97th AMS Annual Meeting, Seattle, WA, USA. 
\item Basu, S. (2016). ``Identifying the limitations of the contemporary planetary boundary layer schemes using an extended self-similarity-based framework'', AMS 22nd Symposium of Boundary Layers and Turbulence, Salt Lake City, UT, USA.
\item DeMarco*, A., and Basu, S. (2016). ``Evidence of extended self-similarity in temperature fields within the mesoscale range'', AMS 22nd Symposium of Boundary Layers and Turbulence, Salt Lake City, UT, USA.
\item He, P., and Basu*, S. (2016). ``Direct numerical simulation of global intermittent turbulence in stably stratified flows'', AMS 22nd Symposium of Boundary Layers and Turbulence, Salt Lake City, UT, USA.
\item Hawbecker*, P., Basu, S., and Manuel, L. (2016). ``Examining microburst scaling relationships through full-physics large eddy simulations'', AMS 22nd Symposium of Boundary Layers and Turbulence, Salt Lake City, UT, USA.
\item Hawbecker*, P., Lu, N.-Y., Basu, S., Kosovi\'{c}, B., and Manuel, L. (2015). ``Coupled mesoscale-LES modeling of a diurnal cycle with a nocturnal low-level jet: the Wangara case study'', June 28 – July 3, AMS 23rd Conference on Numerical Weather Prediction, Chicago, IL, USA.
\item Nunalee, C. G., Basu*, S., Minet, J., and Vorontsov, M. A. (2013). ``Atmospheric refractivity anomalies induced by mesoscale von K\'{a}rm\'{a}n vortex streets'', Imaging and Applied Optics: Propagation Through and Characterization of Distributed Volume Turbulence (pcDVT), OSA, 24-26 June, Washington DC, USA.
\item He*, P., Nunalee, C. G., and Basu, S. (2013). ``Influences of the turbulence parameterizations on atmospheric refractivity simulation and forecasting'', Imaging and Applied Optics: Propagation Through and Characterization of Distributed Volume Turbulence (pcDVT), OSA, 24-26 June, Washington DC, USA.
\item Nunalee*, C. G., and Basu, S. (2012). ``Estimating the higher-order turbulence statistics from LES-generated atmospheric boundary layer flow fields'', AMS 20th Symposium on Boundary Layers and Turbulence, 9-13, July, Boston, MA, USA.
\item Basu*, S., and Holtslag, A. A. M. (2012). ``Why do stable boundary layer simulations over land sometimes crash?'' AMS 20th Symposium on Boundary Layers and Turbulence, 9-13, July, Boston, MA, USA.
\item Kiliyanpilakkil*, V. P., and Basu, S. (2012). ``An algorithm for denoising and gap filling of sodar wind data based on the statistical learning theory'', AMS 20th Symposium on Boundary Layers and Turbulence, 9-13, July, Boston, MA, USA.
\item van de Wiel*, B. J. H., Moene, A., Jonker, H. J. J., Baas, P., Basu, S., Sun, J., and Holtslag, A. A. M. (2012). ``The onset of the very stable boundary layer'', AMS 20th Symposium on Boundary Layers and Turbulence, 9-13, July, Boston, MA, USA.
\item Araya*, G., Castillo, L., Ruiz-Columbi\'{e}, A., Schroeder, J., and Basu, S. (2012). ``On the similarities of the engineering and atmospheric boundary layers'', AMS 20th Symposium on Boundary Layers and Turbulence, 9-13, July, Boston, MA, USA.
\item Basu*, S., and Manuel, L. (2012). ``A large-eddy simulation database for wind energy and boundary-layer meteorology research'', Euromech Colloquium 528 on ``Wind Energy and the impact of turbulence on the conversion process'', 22-24 February, Oldenburg, Germany.
\item Park*, J., Manuel, L., and Basu, S. (2012). ``LES of wind fields and wind turbine load estimation'', 50th AIAA Aerospace Sciences Meeting including the New Horizons Forum and Aerospace Exposition, 9-12 January, Nashville, TN, USA.
\item Park*, J., Manuel, L., and Basu, S. (2012). ``Toward understanding characteristics of the stable boundary layer that influence wind turbine loads'', 50th AIAA Aerospace Sciences Meeting including the New Horizons Forum and Aerospace Exposition, 9-12 January, Nashville, TN, USA.
\item Wang*, Y., Basu, S., and Manuel, L. (2012). ``Realistic turbulence inflow generation for wind turbine design'',  50th AIAA Aerospace Sciences Meeting including the New Horizons Forum and Aerospace Exposition, 9-12 January, Nashville, TN, USA.
\item Basu*, S., Storm, B., and Trier, R. (2010). ``Downscaling the North American Regional Reanalysis wind dataset'', AGU Fall Meeting, 13-17 December, San Francisco, CA, USA. 
\item Basu*, S., Bosveld, F. C., and Holtslag, A. A. M. (2010). ``Stable Boundary Layers with Low-Level Jets: What did we Learn from the LES intercomparison within GABLS?'' 19th Symposium on Boundary Layers and Turbulence, 2-6 August, Keystone, CO, USA.
\item Basu*, S., Bosveld, F. C., and Holtslag, A. A. M. (2010). ``Stable Boundary Layers with Low-Level Jets: What did we Learn from the LES intercomparison within GABLS?'' 5th International Symposium on Computational Wind Engineering, 23-27 May, Chapel Hill, NC, USA.
\item Basu*, S., Bosveld, F. C., and Holtslag, A. A. M. (2009). ``Stable Boundary Layers with Low-Level Jets: What did we Learn from the LES intercomparison within GABLS?'' AGU Fall Meeting, 14-18 December, San Francisco, CA, USA.
\item Gowda*, P. H., Howell, T. A., Hartogensis, O., Basu, S., and Scanlon, B. R. (2009). ``Effect of scintillometer height on structure parameter of the refractive index of air measurements'', AGU Fall Meeting, 14-18 December, San Francisco, CA, USA. 
\item Basu*, S. (2009). ``GABLS 3rd LES intercomparison study'', GABLS Workshop, 26-27 June, Boulder, CO, USA. 
\item Basu*, S., and Holtslag, A. A. M. (2009). ``Solving the mysteries of decoupling, run-away surface cooling, and crashing in stable atmospheric boundary layer model studies'', GABLS Workshop, 26-27 June, Boulder, CO, USA. 
\item Ruiz Columbi\'{e}*, A., Basu, S., and Harshan, S. (2009). ``How does the WRF model capture the intrinsic features of evening transitional boundary layers?'', 10th WRF Users' Workshop, 23-26 June, Boulder, CO, USA. 
\item Harshan, S., Basu*, S., and Ruiz Columbie, A. (2009). ``Evaluating the performance of the WRF model in representing the Antarctic boundary layer'', 10th WRF Users’ Workshop, 23-26 June, Boulder, CO, USA.
\item Sim, C., Basu, S., and Manuel*, L. (2009). ``The influence of stable boundary layer flows on wind turbine fatigue loads'', 47th AIAA Aerospace Sciences Meeting, 5-8 January, Orlando, FL, USA. 
\item Ho, C., Gilliam, X., and Basu*, S. (2009). ``Detecting intermittent turbulence using advanced signal processing techniques'', 47th AIAA Aerospace Sciences Meeting, 5-8 January, Orlando, FL, USA. 
\item Basu*, S., Holtslag, A. A. M., and Bosveld, F. (2008). ``Implications of the GABLS LES intercomparison studies for future wind energy projects'', AGU Fall Meeting, 15-19 December, San Francisco, CA, USA.
\item Basu*, S., Holtslag, A. A. M., van de Wiel, B. J. H., Moene, A. F., and Steeneveld, G. J. (2008). ``Fundamental limitations of using sensible heat flux as a lower boundary condition in stable boundary layer simulation and modeling'', 18th Symposium of Boundary Layers and Turbulence, 9-13 June, Stockholm, Sweden.
\item Basu*, S., Steeneveld, G. J., Holtslag, A. A. M., and Bosveld, F. C. (2008). ``Large-eddy simulation intercomparison case setup for GABLS3'', 18th Symposium of Boundary Layers and Turbulence, 9-13 June, Stockholm, Sweden.
\item Phillipson*, J. A., Basu, S., and Gilliam, X. (2008). ``Identifying differences in polar and mid-latitude turbulence bursting events'', 88th AMS Annual Meeting, 20-24 January, New Orleans, LA, USA.
\item Anderson*, W. C., Basu, S., and Letchford, C. W. (2007). ``Next-generation subgrid-scale closures for large-eddy simulation of turbulent flow in the atmospheric boundary layer'', 12th International Conference on Wind Engineering, 1-6 July, Cairns, Australia.
\item Storm*, B., Dudhia, J., and Basu, S. (2007). ``WRF PBL sensitivities and verification in forecasting low-level jet events.'' 8th Annual WRF user’s workshop, 11-15 June, Boulder, CO, USA.
\item Basu*, S., and Vinuesa, J.-F. (2007). ``Dynamic LES modeling of observed diurnal cycles.'' AGU Joint Assembly, 22-25 May, Acapulco, Mexico.
\item Foufoula-Georgiou*, E. and Basu, S. (2007). ``Advances in precipitation forecast verification: The forecast quality index and a case study in assessing WRF model performance.'' AGU Joint Assembly, 22-25 May, Acapulco, Mexico.
\item Basu*, S., Foufoula-Georgiou, E., Lashermes, B. and Arnéodo, A. (2007). ``Estimating intermittency exponent in neutrally stratified atmospheric surface layer flows: A robust framework based on magnitude cumulant and surrogate analyses.'' AGU Joint Assembly, 22-25 May, Acapulco, Mexico.
\item Basu*, S., and Foufoula-Georgiou, E. (2007). ``Spatial forecast verification methods intercomparison project: Forecast Quality Index'', Verification Method Intercomparison Meeting, 19-24 February, Boulder, CO, USA. 
\item Basu*, S., and Vinuesa, J.-F. (2006). ``Dynamic LES modeling of a diurnal cycle.'' AGU Fall Meeting, 11-15 December, San Francisco, CA, USA. 
\item Basu*, S., Lashermes, B., Foufoula-Georgiou, E., and Arnéodo, A. (2006). ``Does atmospheric stability affect inertial-range properties of boundary layer turbulence?'' AGU Fall Meeting, 11-15 December, San Francisco, CA, USA. 
\item Basu*, S., Anderson, W., and Vinuesa, J.-F. (2006). ``Dynamic LES modeling of diurnal cycles.'' AMS 17th Symposium on Boundary Layers and Turbulence, 22-25 May, San Diego, CA, USA.
\item Gilliam, X., and Basu*, S. (2006). ``Identifying intermittent structures in stable boundary layer flows.'' AMS 17th Symposium on Boundary Layers and Turbulence, 22-25 May, San Diego, CA, USA.
\item Basu*, S. (2006). ``Large-eddy simulation of stable boundary layers: a case study.'' 1st iLEAPS Science Conference, 21-26 January, Boulder, CO, USA.
\item Mehta, K., Peterson, R., Basu*, S., and Swift, A. (2006). ``Doctoral degree in Wind Science and Engineering'', 3rd National Conference on Wind Engineering, 5-7 January, Calcutta, India.
\item Basu*, S. (2006). ``Large-eddy simulation of stable boundary layer turbulence.'' 3rd National Conference on Wind Engineering, 5-7 January, Calcutta, India. 
\item Porté-Agel*, F., Foufoula-Georgiou, E., Stoll, J. R., Basu, S., and Venugopal, V. (2005). ``Surface heterogeneity effects on land-atmosphere fluxes from tuning-free large-eddy simulations.'' 5th International Scientific Conference on the Global Energy and Water Cycle, 20-24 June, Orange County, CA, USA.
\item Basu*, S., Porté-Agel, F., and Foufoula-Georgiou, E. (2004). ``Development of stable boundary layer parameterizations in numerical weather prediction models: a scale-dependent dynamic LES modeling approach.'' AMS 16th Symposium on Boundary Layers and Turbulence, 9-13 August, Portland, ME, USA.
\item Basu*, S., Porté-Agel, F., and Foufoula-Georgiou, E. (2004). ``Revisiting a few outstanding issues related to the stably stratified atmospheric boundary layer turbulence research.'' European Geosciences Union General Assembly, 25-30 April, Nice, France.
\item Basu*, S., and Foufoula-Georgiou, E. (2001). ``Information theoretic approach to questions of precipitation predictability.'' AGU Fall Meeting, 10-14 December, San Francisco, CA, USA.
\item Basu*, S., and Foufoula-Georgiou, E. (2001). ``A framework for the detection of nonlinearity and chaoticity: application to high frequency storm rates.'' 7th International Conference on Precipitation, 30 June - 3 July, Rockport, ME, USA. 
\item Basu*, S., Henshaw, P. F., and Biswas, N. (2000). ``Perspectives of adsorption in the new millennium.'' Canadian Society for Civil Engineering, 6th Environmental Engineering Specialty Conference, 8-10 June, London, Canada, pp. 408-413. 
\item Basu*, S., and Henshaw, P. F. (2000). ``Correlations of relative adsorbability with physical and structural properties of VOCs -- application of an artificial neural network.'' Canadian Society for Civil Engineering, 6th Environmental Engineering Specialty Conference, 8-10 June, London, Canada, student competition paper. 
\end{etaremune}

